% THETA3-B20 contribution by /root/definite_bridge, 21 November 2026.
\subsection{THETA3-B20: isotropic column completion preserves a first determinant value}\label{B20}

For the exponent-one determinant occurring in the original
canonical wedge, the extra section columns required by
\texttt{IsCols} can be supplied without changing the value of
the differential at a fixed CM point.  The reference theta-null
values and the integral formal CM action must be retained.
This is a finite-jet comparison; it does not identify the
resulting theta sections as global automorphic functions.

\paragraph{The elementary determinant identity.}
Let $G$ be a commutative formal group over a characteristic-zero
field, with commuting invariant derivatives $D_1,\ldots,D_i$.
For a section germ use a fixed normalized canonical frame, and
write its scalar coefficient as $F$.
Let $F_1,\ldots,F_i,R_1,\ldots,R_i$ be scalar formal series with
\[
                  R_c(0)=r_c,\qquad D_aR_c(0)=0
                     \quad(1\leq a,c\leq i).
\]
On $2i$ copies write $D_{a,c}$ for $D_a$ on the $c$-th copy.
For any scalar $\beta$ define
\[
 \mathcal D_\beta=
     \det(D_{a,c}+\beta D_{a,c+i})_{1\leq a,c\leq i},
 \qquad
 \mathcal D_0=\det(D_{a,c})_{1\leq a,c\leq i}.
\]
Then
\begin{equation}\label{B20-value}
 \mathcal D_\beta
       \left(\prod_{c=1}^iF_c(t_c)R_c(t_{c+i})\right)\Big|_0
     =\left(\prod_{c=1}^ir_c\right)
          \mathcal D_0\left(\prod_{c=1}^iF_c(t_c)\right)\Big|_0.
 \tag{B20a}
\end{equation}
Expand the determinant by permutations and each of its $i$
factors into its two summands.  Every term selecting an extra
column differentiates its particular $R_c$ exactly once;
that term vanishes at zero.  The only surviving choice selects
all the original columns and leaves the reference factors
undifferentiated.  This proves (B20a), with its sign and without
a factorial or a conjugation of $\beta$.

\paragraph{Actual integral even references.}
Let $R$ be a complete discrete valuation ring of residue
characteristic $p\ne2$, and let $A/R$ be an abelian scheme
with a symmetric relatively ample line bundle $L$, rigidified
at zero.  Normalize the symmetry on $L^k$ to be the identity
on its fibre at zero.  Suppose the evaluation map
\[
                  H^0(A,L^k)\longrightarrow (L^k)_0=R
 \tag{B20b}
\]
is surjective.  This is an explicit condition on the chosen
theta level.  It holds for sufficiently large $k$, by relative
generation and cohomology-and-base-change for ample powers.
It can be imposed by choosing a sufficiently large square
level prime to $p$; any pullback isogeny and its derivative
factor must be recorded when an already fixed input is moved
to that level.

Choose an integral section $s$ with $s(0)=1$ and put
\[
                        r=\frac{s+[-1]^*s}{2},
\]
using the normalized symmetry to regard both summands as
sections of the same line bundle.  Then $r(0)=1$ and its
first derivative in any symmetric origin frame is zero,
because $d[-1]=-1$.  The canonical potential of A8--A9 has
zero constant and linear terms: it is obtained from a
biextension series vanishing on both axes by restriction
to $(x,[w_0]x)$.  The actual canonical-frame first jets of
$r$ therefore vanish as well.  Thus $r$ supplies the reference
germ in (B20a), with $r_c=1$ and with an integral section,
without assuming the integrality of its entire canonical
theta-frame expansion.

If the chosen reference has a unit value $r_c$ rather than one,
division by that unit gives the normalized reference.  If its
value is merely nonzero in the fraction field, the value can
still be divided out in characteristic zero, but that operation
has a valuation cost.  An assertion of integral divisibility
must then retain that cost.  Merely choosing a nonvanishing
complex theta-null value does not prove (B20b).

\paragraph{The actual completed CM columns.}
Assume the formal group has an integral action of $\OO_E$,
all its endomorphisms are defined over $R$, and its coefficientwise
continuous extension to $\OO_E\otimes\Z_p$ is available.
This extension can also be constructed directly: use the formal
logarithm over the fraction field, write an endomorphism as
$\operatorname{Exp}\circ d[a]\circ\log$, and take the
coefficientwise limit of integral global CM actions.  Each finite
jet is a polynomial over the fraction field in the entries of
$d[a]$, and closedness of $R$ preserves integrality.

Fix one CM embedding direction $\tau$ in which the tangent
frame $u_1,\ldots,u_i$ is integral and saturated.  Choose
\[
 b\in(\OO_E\otimes\Z_p)^\times,
                   \qquad b\overline b=-1,
                   \qquad \beta=d[b]|_\tau\in R^\times.
 \tag{B20c}
\]
The bar in (B20c) is the CM involution on the completed algebra;
it is not an assertion about the ordinary complex conjugate
of an arbitrarily chosen complex image of $\beta$.

Such a $b$ exists when $p$ is odd and unramified in $E$.
At a split pair in $E/F$, take the two unit components
$(1,-1)$.  At an unramified quadratic local extension with
residue-field sizes $q^2/q$, the Teichm\"uller units contain an
element with $b^{q+1}=-1$: $2(q+1)$ divides $q^2-1$ because
$q$ is odd.  Its local CM norm is $-1$.
Make these choices at every factor of $E\otimes\Q_p$.
The embedding values required in (B20c) are roots of unity or
$\pm1$, so they can be included in the coefficient field after
a finite extension.

The matrix of section columns is
\[
                         B=[I_i\mid\beta I_i].
 \tag{B20d}
\]
It has an $i\times i$ unit minor.  For every $N$, choose
$a_N\in\OO_E$ congruent to $b$ modulo
$p^N(\OO_E\otimes\Z_p)$.  Then
\[
 X_N=[I_i\mid a_NI_i],\qquad
 X_N\overline{X_N}^{\,t}
                  =(1+a_N\overline{a_N})I_i
                        \equiv0\pmod{p^N\OO_E}.
\]
At the selected embedding $X_N$ converges to $B$.
Thus (B20d) satisfies the literal supported-block, unit-minor
and CM-isotropic-approximation requirements of
\texttt{IsIsoCols}; its single block has precisely $2i$ columns.
For several determinant factors the construction uses disjoint
blocks, as required by \texttt{IsCols}.  The precise definitions
are
\texttt{Development/AbelianVarietyTheory/Definitions/Divisors/NonemptyValExtOfOverLa.lean:331--342},
read in A9.  The differential polynomial produced by this
matrix is exactly $\mathcal D_\beta$, since the selected
embedding acts on every $u_a$ by the same scalar $\beta$.

Applying A9's integral-unit lemma to these columns makes the
entire canonical theta-frame factor an integral unit after
the corresponding formal substitution.  Applying (B20a) to
the actual section germs then identifies the exponent-one
value with the original determinant value, including the
reference theta-null factors.  Any use of a particular
\texttt{nuFac} lower bound still requires its stated local
type, saturation and valuation-normalization hypotheses.

\paragraph{The two different oscillator tori in B19.}
The construction can be made on one common polarized abelian
variety without identifying $B(d_1)$ with $B(d_2)$.
Let $A_1,A_2$ be the two actual CM theta tori at the fixed
rational negative line, with line bundles $M_1,M_2$ at the
chosen common input levels.  Choose even reference sections
$r_i$ of $M_i$ with $r_i(0)=1$ under (B20b), and put
\[
 A=A_1\times A_2,\qquad M=M_1\boxtimes M_2,
 \qquad S_1=s_1\boxtimes r_2,
 \quad S_2=r_1\boxtimes s_2,
 \quad R_0=r_1\boxtimes r_2.
 \tag{B20e}
\]
All three are actual sections of the same ample line bundle
$M$; $R_0(0)=1$ and its first canonical jets vanish.
For the tangent rows take
\[
             u_a=(u_{1a},u_{2a})\in
                          \operatorname{Lie}(A_1)\oplus
                          \operatorname{Lie}(A_2),\quad a=1,2.
\]
The $u_{ia}$ are the normalized algebraic $\mu$ directions
from B19, with its period factors kept outside the algebraic
jet value.  The diagonal $E$-action has the same eigenvalue
on both summands in this selected embedding direction.
Under the stated integral saturated-frame hypothesis they
therefore admit the columns (B20d).

The product rule and $dr_i(0)=0$ give
\[
       D_{u_a}S_1(0)=D_{u_{1a}}s_1(0),\qquad
       D_{u_a}S_2(0)=D_{u_{2a}}s_2(0).
\]
Consequently the original two-torus first-jet determinant is
the two-column determinant for $S_1,S_2$ on $A^2$.
By (B20a) it is also the four-column isotropic determinant
for $S_1,S_2,R_0,R_0$ on $A^4$.  Nonunit changes of tangent
frame or of level must be kept as explicit factors; the
saturation condition is not supplied by an arbitrary
complex choice of frame.

This works simultaneously for the finite list of transformed
sections in (B19j): they lie in the same finite section spaces
on the fixed $A_i$, and the references can be chosen once.
It therefore completes the exponent-one \emph{value} comparison
for that finite CM-jet functional.  It does not assert that
$A$ has the particular uniform CM Lie multiplicities required
by every recovered \texttt{Datum}, or supply a missing
integral model, tame-generator normalization, or global
\texttt{MainX} witness.

\paragraph{Why the conclusion stops at exponent one.}
For $\mathcal D_\beta^{a}$ with $a>1$, a reference factor
may receive two or more derivatives, and evenness no longer
kills all the added terms.  A sufficient replacement is that
each reference have constant canonical jet through order $a$:
then every extra derivative term still vanishes.  Jet separation
over $\overline\Q$ alone does not give this with integral unit
normalization.  The inverse canonical gauge can have
nonintegral higher coefficients, as C22 demonstrates.
No such high-order integral references are asserted here.

Even at exponent one, the completion is not an equality of
the original automorphic functions on the ball.  Away from
the fixed CM point the extra even theta sections contribute
their varying theta-null factors and additional line-bundle
weights.  Their chosen value at the point is all that was
used in (B20a).  Thus B20 removes a specific finite-jet column
obstruction while leaving the global section and high-exponent
comparisons explicit.\hfill$\square$
